
\section{Introduction and Related Works}
The fundamental promise of transformation-based Privacy-Preserving Face Recognition (PPFR) systems is compelling: verify a user's identity without ever exposing their raw facial data \cite{dai_fracface_2025, mi_privacy-preserving_2023, mi_privacy-preserving_2024, mi_duetface_2022, jin_faceobfuscator_2024, wang_privacy-preserving_2022}. A user's face is converted into a protected template that a recognition server can use for matching, but that, ideally, cannot be reversed to recover the original image or used to infer private attributes. The primary adversary in this setting is the curious or malicious service provider who receives the template \citep{erkin_privacy-preserving_2009, ji_privacy-preserving_2022}.

The prevailing evaluation paradigm for these systems is currently significantly limited. Historically, the dominant objective has been to prevent the reconstruction of the original registration image, with attack success measured through pixel-level or local similarity metrics such as peak signal-to-noise ratio (PSNR) and structural similarity index measure (SSIM), a legacy inherited from image privacy and compression literature. While some recent approaches, such as FracFace \cite{dai_fracface_2025}, also report identity similarity and employ StyleGAN for non–pixel-level reconstruction, their still have a strong reconstruction-centric focus. A large body of prior work adopts this paradigm, using resistance to pixel-level recovery as evidence of privacy protection and optimizing attack objectives accordingly. Representative systems, including DuetFace \cite{mi_duetface_2022}, MinusFace \cite{mi_privacy-preserving_2023}, PartialFace \cite{mi_privacy-preserving_2024}, FaceObfuscator \cite{jin_faceobfuscator_2024}, and FracFace, explicitly rely on this reconstruction-based evaluation to argue robustness against recovery attacks.

\begin{figure}
    \centering
    \includegraphics[width=\linewidth]{imgs/vis.pdf}
    \caption{Regeneration attack results. In each subplot, the left is the original image and the right is the regenerated image from the protected template. Each row shows examples from one PPFR method, in the order of PartialFace, MinusFace, and FracFace. }
    \label{fig:vis}
\end{figure}

This evaluation paradigm, however, rests on a critical implicit assumption: that preventing pixel-level reconstruction is both necessary and sufficient to prevent identity leakage. In this paper, we show that this assumption does not hold. Crucially, compromising privacy does not require recovering the original registered image, nor does pixel-level similarity reliably correspond to identity consistency. In the facial domain, two images that are visually or pixel-wise similar may represent different identities and images that represent same identity could have completely no pixel or local similarity. CanFG \cite{wang_make_2024} can generate two images with very high pixel-level similarity yet in completely different identities; conversely, in daily life, any two arbitrary photos of the same person (one ID photo and one social media photo) would have very high identity similarity yet likely very low pixel-level or structure-level similarity. An example is provided in Figure~\ref{fig:canfg}. 
%; furthermore, we demonstrate in Appendix~\ref{sec:bypass} that generating high-fidelity images that trivially bypass pixel-level metrics without preserving identity is a computationally negligible task.

This misconception mislead not only the evaluation but also the simulated attack design because the overlooked fact that identity-revealing information can remain accessible even when pixel-level reconstruction is infeasible. By employing pixel-level loss functions, simulated attackers (red-team researchers) are inadvertently trapped into pursuing the specific registration image as ground truth. This objective is often mathematically impossible due to the information loss in protection, causing the generator to produce a blurry image, likely an average of all images in the dataset. The failure is illustrated in Figure~\ref{fig:vsgan}. Although FracFace \cite{dai_fracface_2025} also reports identity similarity, the simulated U-Net attack remains reconstruction-based, aiming to recover the original pixels. The StyleGAN \cite{karras_style-based_2019} attack also failed to generate a face with the same identity as the original image, as shown in Figure~\ref{fig:vsgan}.

\begin{figure}
    \centering
    \begin{minipage}[t]{0.4\linewidth}
        \centering
        \includegraphics[width=\linewidth]{imgs/canfg.png}
        \captionof{figure}{SSIM and PSNR are not always correlated with identity correlation.}
        \label{fig:canfg}
    \end{minipage}
    \hfill
    \begin{minipage}[t]{0.4\linewidth}
        \centering
        \includegraphics[width=\linewidth]{imgs/vsgan.png}
        \captionof{figure}{Using pixel-level loss or StyleGAN will yield unsuccessful reconstruction compared to our ID-guided method.}
        \label{fig:vsgan}
    \end{minipage}
\end{figure}

\begin{table}
    \centering
        \caption{Comparison between pixel-level metrics and identity-level metrics on two cases: a protected face generated by CanFG with its original face, and two images of the same person. We can see that higher pixel-level similarity does not mean higher identity-level similarity. }
    \label{tab:canfg}
    
    \begin{tabular}{c|cccc}
    \toprule
         Compared With &  SSIM&  PSNR&  MSE& FS\\
         \midrule
         Another Face&  0.235&  10.44&  6699& 0.586\\
         CanFG Face&  0.841&  26.81&  143& 0.008\\
         \bottomrule
    \end{tabular}
    %\vskip -0.1in
\end{table}

Recent work has begun to challenge the reconstruction-centric perspective in non-facial domains by shifting attention toward semantic-level inversion. In such settings, the attacker aims not to recover the original image itself, but to regenerate information that is semantically consistent with the original identity. This goal is both easier to achieve and more aligned with realistic attack objectives. For instance, \citet{yue_gradient_2023} proposes a semantic recovery framework that leverages generative models to synthesize semantically consistent images without pixel-level similarity in federated learning attacks. Complementary findings further suggest that traditional reconstruction metrics fail to capture how humans perceive privacy leakage~\cite{sun_privacy_2023}. 

Our approach differs from standard Model Inversion Attacks (MIAs). One category of MIA, represented by \citet{wang_diffusion-driven_2025}, recovers original images from embeddings. For many deep facial embedding models, such processes resemble image generation tasks rather than adversarial attacks, as they exploit the inherent invertibility of learned representations. Similar works exist in the ID-controlled image generation domain, like Arc2Face \cite{papantoniou_arc2face_2024}, FulID \cite{guo_pulid_2024}, and FaceID IP-Adapter \cite{ye_ip-adapter_2023}. In contrast, our method targets structural vulnerabilities in the template generation process itself. Since this conversion often utilizes rule-based transformations independent of specific deep models, the attack surface differs from embedding-based reconstruction. Another category of MIA, which is closer to the original definition \cite{fredrikson_model_2015}, such as the ones prevented by \citet{wang_adaptive_2024}, aims to reconstruct the training dataset to compromise identity privacy in \textit{the training set}. This focus on training data deviates from our objective of protecting individual user templates. Existing solutions for training-level privacy include FaceMAE \cite{wang_facemae_2022} and the use of synthetic data \cite{bae_digiface-1m_2022}. 





\textbf{Our Contribution:} This paper argues that the prevailing evaluation paradigm for frequency-domain PPFR, which measures privacy as resistance to pixel-level reconstruction, does not capture identity-level leakage. We introduce and validate an identity-centric evaluation standard, and present FaceLinkGen, an attack that performs identity linkage and face regeneration directly from protected templates without recovering original pixels. We show that pixel-level metrics can overestimate the privacy provided by frequency-domain PPFR systems: identity information can be extracted via a straightforward distillation pipeline even when the original image remains ``unreconstructable'' under PSNR and SSIM. The simplicity of the pipeline is intentional: if a standard procedure suffices, the vulnerability is in the representation. This work encourages identity-centric evaluation alongside existing metrics in future PPFR research.






\section{Threat Model}
The PPFR paradigm was originally motivated by a specific and well-defined concern: a user must submit face data to a remote recognition server they do not fully trust \citep{erkin_privacy-preserving_2009}. Since the raw face image is deeply personal, the user has reason to want the server to perform recognition without gaining access to the original biometric. PPFR addresses this by transforming the face into a protected representation before transmission, one that the server can use for matching but, ideally, cannot revert to recover the original face, infer soft biometrics such as age, gender, or ethnicity, or link the identity to records outside the service. The primary adversary in this setting is therefore the \emph{curious or malicious service provider}, an insider who legitimately receives the protected template and seeks to exploit it.

Protection against network-level eavesdroppers is a separate concern and was not the original intent of PPFR. Interception of data in transit is well addressed by general-purpose secure communication protocols such as TLS, which are application-agnostic and do not require altering the face representation itself. Nevertheless, some recent PPFR work has adopted an external-attacker framing: \citet{mi_privacy-preserving_2024} describes this attacker as ``typically envisioned as a malicious third-party wiretapping the transmission,'' departing from the original insider-centric design intent.

We return to first principles and evaluate PPFR against its intended adversary: a service provider with oracle access to the conversion process (the ability to query it with arbitrary inputs and observe outputs), but no knowledge of its internal architecture or parameters. Even under this insider-focused framing, our attacker model assumes less than prior work: \citet{mi_privacy-preserving_2024} assumes knowledge of the conversion architecture (but not the random channel selection parameters), and \citet{mi_privacy-preserving_2023} assumes access to the conversion process but not the selected channel IDs. We assume no knowledge of the architecture, parameters, or hyperparameters whatsoever.

For completeness, we also note that \citet{zhang_validating_2024} requires approximately 6{,}900 online verification queries per identity and depends on the server returning continuous similarity scores for optimization, a setup highly susceptible to rate limiting and fraud detection. Many deployed FR systems return only binary accept/reject decisions or quantized similarity scores, rendering such approaches infeasible \citep{noauthor_characterizing_2011, noauthor_bioapi_2001}. Our method does not rely on server-side behavior at all.











\section{Methods}
The simplicity of our method is intentional. We show that even strong protection methods fail with a simple, standard distillation process, proving that the vulnerability resides in the representation itself. To formulate this, we consider a face image $X$ as a combination of identity information $z_I$ and non-identity (nuisance) information $z_N$, such that $X \sim p(\cdot \mid z_I, z_N)$.

In transformation-based PPFR systems, a protected template $T$ is generated to hide the visual data of $X$ while retaining identity utility. This process can be viewed as a lossy mapping that suppresses the information quantity of $z_N$ while preserving $z_I$:
\begin{equation}
T \sim p(\cdot \mid z_I).
\end{equation}
Existing evaluations often equate privacy with the failure of pixel-level reconstruction. However, since $z_N$ is largely discarded, recovering the original pixels $X$ is a severely ill-posed problem. Conventional attacks fail because they attempt to optimize for specific nuisance factors (e.g., exact lighting or pose) that no longer exist in $T$, resulting in blurry or identity-inconsistent outputs.

Our approach, FaceLinkGen, instead focuses on extracting the remaining identity information. We use a distillation-style procedure to align the template domain with a standard identity embedding space. Given a public dataset, we train a student model $f_s$ to recover an identity representation $z_I'$ from $T$. The training objective is to maximize the cosine similarity between the student's output and the embeddings $z_t$ extracted by a frozen teacher model $f_t$ from original images $X$:
\begin{equation}
\mathcal{L} = 1 - \frac{1}{N}\sum_{k=1}^{N} s\bigl(f_s(t_k), f_t(i_k)\bigr),
\end{equation}
where $z_I' = f_s(T)$ is the identity feature recovered by the attacker. 

Once $z_I'$ is extracted, the attack bypasses the need for pixel reconstruction by leveraging a diffusion-based generative model $g_{\mathrm{diff}}$. Rather than trying to find the original $z_N$, we substitute the missing information by sampling from the model's stochastic noise $\epsilon$:
\begin{equation}
\epsilon \sim \mathcal{N}(0, I), \quad
Y = g_{\mathrm{diff}}(z_I', \epsilon).
\end{equation}
In this process, we do not attempt to recover the discarded original nuisance factor $z_N$. Instead, we extract the identity representation $z_I' = f_s(T)$ from the template and introduce a stochastic noise vector $\epsilon \sim \mathcal{N}(0, I)$ as a proxy for a newly sampled set of non-identity factors $z_N'$. This enables the model to bypass the ill-posed reconstruction of the original $z_N$ by providing the necessary information density to synthesize a realistic face $Y$ with the same identity but different attributes. ($z_N \not\approx z_N'$) Our formulation demonstrates that as long as the information quantity of $z_I$ persists in $T$, an attacker can recover $z_I'$ and combine it with a random $\epsilon$ to regenerate a high-fidelity, identity-consistent face. This implies that, for frequency-domain PPFR methods, visual distortion of $z_N$ does not prevent identity extraction, as recovering $z_N$ is not a prerequisite for a successful attack.

Note that our attack method is independent of the specific face recognition model employed by the PPFR server. We use ArcFace as the teacher and student network solely because it is a widely-adopted, publicly available facial embedding model, and its compatibility with the Arc2Face generative model simplifies our demonstration. Other embedding models like FaceNet \cite{schroff_facenet_2015} can also be used as long as there is a compatible generative model. The server's actual recognition backbone could be any commercial or proprietary model. Our attack only requires that the protected template retains identity-discriminative features that are learnably aligned with some facial embedding space accessible to the attacker, a condition implicitly assumed by any PPFR system that aims to preserve recognition utility.

\section{Attack Vectors}
The identity extraction framework described above enables two distinct attack vectors. The linkage attack requires only the extracted embedding $z_I'$ and operates entirely in the embedding space, while the regeneration attack additionally employs the diffusion-based generator $g_{\mathrm{diff}}$ to synthesize identity-consistent face images.
\subsection{Linkage Attack}
A linkage attack aims to associate a real-world identity (e.g., a public face image) with a protected identity, or to link two protected templates belonging to the same individual across different leaked databases. The first case is referred to as face-to-template linkage, while the second is template-to-template linkage. ISO/IEC 24745 explicitly requires resistance against template-to-template search, but does not address face-to-template search. This is likely a utlity trade off for the verification needs of the service provider. This attack vector is similar to an attack vector for hashing: when the input space is known, an attacker can enumerate all possible inputs and map each hashed output back to its original input.


In both attack scenarios, the attacker first obtains a query embedding $e_q$. This embedding can be extracted using either the student model $f_s$ or the teacher model $f_t$, depending on the domain of the query data. The adversary then computes embeddings for all protected templates in the leaked database using $f_s$ and performs a nearest-neighbor search. This process can be written as
\begin{equation}
\arg\max_{t \in T} \; s(e_q, f_s(t)),
\end{equation}
where $s(\cdot, \cdot)$ denotes cosine similarity.


\begin{figure*}
    \centering
    \includegraphics[width=\linewidth]{imgs/distill_2.png}
    \caption{After our model is trained, two main attack vectors can be performed: linkage and re-generation.}
    \label{fig:placeholder}
\end{figure*}
\subsection{Regeneration Attack}

As discussed earlier, reconstructing the original enrollment image is unnecessary and, in many cases, impossible. However, once a universal face embedding (e.g., ArcFace) can be extracted from a protected template, modern face generation models can be leveraged. In this work, we use Arc2Face, which takes a facial embedding as input and generates a face image whose embedding matches the input. This allows us to synthesize a realistic face corresponding to the protected template without reconstructing the original image.

















\section{Experiments and Results}
We selected three frequency-domain PPFR methods with accessible source code: PartialFace~\cite{mi_privacy-preserving_2023} (ICCV 2023), MinusFace~\cite{mi_privacy-preserving_2024} (CVPR 2024), and FracFace~\cite{dai_fracface_2025} (NeurIPS 2025). These represent the mainstream of currently open-sourced transformation-based PPFR work. For distillation, we used a subset of CASIA-WebFace~\cite{yi_learning_2014} with around 10K identities and 90K images. Importantly, there is no dataset or architecture overlap between the attacked methods and our student model: PartialFace and FracFace are training-free hardcoded transformations, and MinusFace is trained on MS1Mv2~\cite{guo_ms-celeb-1m_2016}, while our attack model is trained on CASIA-WebFace. The facial embedding model is Antelopev2 with one additional 3x3 Conv2D layer added before to be compatible with different template formats (channel numbers) if needed. The Antelopev2 is chosen because it is what Arc2Face accepts. The dataset is split into a training and a validation set in an 80-20 ratio. For regeneration testing, we used three datasets: the validation hold-out set of CASIA-WebFace, ``this person does not exist'' (TPDNE) dataset \footnote{TLeonidas/this-person-does-not-exist}, and Labelled Faces in the Wild (LFW) dataset. The hold-out set is used to test the ability of our method in real images while ensuring no ID duplications. The LFW dataset is used to test the cross-dataset performance of our method with distribution shift, and the TPDNE is used as a synthetic dataset to avoid data cross-contamination from Stable Diffusion 1.5 and Arc2Face training data. Compared to the hold-out set and LFW, photos in the TPDNE dataset are also closer to a headshot, which is what is usually used to create the protected templates.

The distillation process was completed in under two hours on a single NVIDIA A6000 GPU for each of the three evaluated methods, at an estimated cost of approximately USD 0.80 to 1.60. The same process can also be executed on consumer GPUs such as the RTX 4090 or RTX 5090 with comparable wall-clock time; in fact, it is theoretically faster on the RTX 5090 due to the newer architecture of the RTX 5090. These extreme low costs are deliberate: they serve as a lower-bound analysis demonstrating that current protection mechanisms succumb to a lightweight, generic distillation without requiring complex adversarial optimization.


\subsection{Minimal-Resource Attack}
To stress-test the data requirements, we conducted an experiment on FracFace with only approximately 800 images from 100 identities. With higher weight decay and lower batch size to mitigate overfitting, training completed in under 50 seconds, yet the attack still achieved a 97.0\% generation pass rate at FAR $1\times10^{-5}$ and 99.5\% linkage accuracy. We further reduced the training set to 256 images from 32 identities, obtaining 98.7\% linkage accuracy and 90.5\% regeneration success. The consistently high performance across these resource-constrained settings suggests that identity information is not meaningfully disrupted by the transformation, but rather preserved in a form that remains easily extractable.

\subsection{Linkage Attack}

Since all templates are converted to the standard ArcFace domain, we can not only link between original images and protected templates, but also link between two templates from the same or different protection methods. In any case, we are linking two different images (or templates) of the same person, not a face image and its corresponding template. We used the CASIA-WebFace hold-out set for the linkage attack to ensure no identity overlap between training and testing identities. The hold-out dataset size is 2115.  

\begin{table}
    \caption{Comparison of protection claims. FracFace measures the proportion of distorted frequency channels; our metric measures identity recovery success via commercial verification. We demonstrate that high channel disruption, as reported in prior work, does not prevent identity extraction.}
    % %\vskip 0.15in
    \label{tab:res}
    
    \centering
    \scalebox{1.0}{
    \begin{tabular}{ccccc}
    \toprule
         &  Success@5&  Pass@1e-5& Pass@1e-4 &Pass@1e-3 \\
         \midrule
         \multicolumn{5}{l}{\textit{Dataset: TPDNE}} \\
         \midrule
         PartialFace&  1.000&  0.993& 0.996&0.998\\
         MinusFace&  0.996&  0.936& 0.970&0.989\\
         FracFace&  0.992&  0.904& 0.957&0.985\\
      \midrule
         \multicolumn{5}{l}{\textit{Dataset: CASIA-WebFace Hold-Out}} \\
         \midrule
         PartialFace&  0.992&  0.957& 0.970& 0.982\\
         MinusFace&  0.989&  0.930& 0.958& 0.978\\
         FracFace&  0.991&  0.920& 0.950& 0.977\\
               \midrule
         \multicolumn{5}{l}{\textit{Dataset: LFW}} \\
         \midrule
         PartialFace&  0.988&  0.980& 0.983& 0.986\\
         MinusFace&  0.987&  0.974& 0.981& 0.983\\
         FracFace&  0.979&  0.943& 0.961& 0.970\\
         % Another Image&   &  & & \\
         \bottomrule
    \end{tabular}}
    % \vspace{-0.1in}
\end{table}
The closed-set 1-to-N linkage results are in Table~\ref{tab:linkage}. The original-image-to-original-image linkage (0.88) establishes a performance upper bound. With the WebFace dataset containing 9.3\%-13.0\% noise \cite{wang_devil_2018}, perfect linkage is impossible regardless of method. Our attack achieves linkage success rates consistently above 70\%, frequently exceeding 80\%, essentially reaching the dataset's theoretical maximum performance. This confirms that the extracted embeddings function as effective identity descriptors for cross-domain matching. Additionally, the 1-to-1 verification accuracy used in the traditional face recognition benchmark (Table~\ref{tab:1v1}) remains near 100\% and comparable to the original ArcFace performance, demonstrating that the protection systems fail to meaningfully impede identity matching.


\begin{table}[]
    \centering
    \caption{1-to-1 verification accuracy between template-to-face and face-to-face on LFW.}
    \label{tab:1v1}
    \begin{tabular}{ccc}
    \toprule
         &  Accuracy& AUROC\\
         \midrule
         MinusFace-to-Face&  0.992& 0.995\\
         FracFace-to-Face&  0.988& 0.993\\
         PartialFace-to-Face&  0.992& 0.996\\
         Face-to-Face&  0.998& 0.998\\
         \bottomrule
    \end{tabular}
    % \vspace{-0.1in}
\end{table}

\begin{table}
\centering
\caption{Linkage Results between MinusFace, PartialFace, FracFace, and Original Image Embeddings on CASIA-WebFace dataset. The numbers reported are top-1 recall at a closed set setting.}
\label{tab:linkage}

\begin{tabular}{lrrrr}
\toprule
\text{Query} & \multicolumn{4}{c@{}}{\text{Key}}\\
    \cmidrule(l){2-5}
  & FracFace  & MinusFace  & PartialFace  & Original\\
\midrule
FracFace  & 0.7863 & 0.7537 & 0.8137 & 0.8478 \\
MinusFace  & 0.7305 & 0.7206 & 0.7754 & 0.8132 \\
PartialFace  & 0.8028 & 0.7868 & 0.8270 & 0.8572 \\
Original& 0.8444 & 0.8241 & 0.8563 & 0.8823 \\
\bottomrule
\end{tabular}
\end{table}

\begin{table}
    \centering
    \caption{Cross-validation results (pass rate) using the Amazon API }
    \label{tab:amazon}
    
    \begin{tabular}{ccc}
    \toprule
         PartialFace&   MinusFace&  FracFace\\
         \midrule
        0.99 &   0.98&  0.92\\
         \bottomrule
    \end{tabular}
    % \vspace{-0.1in}
\end{table}
\subsection{Regeneration Attack}

For the regeneration attack, we evaluate identity recovery on the first 1,000 images from each of the following datasets: the TPDNE dataset, the hold-out set of CASIA-WebFace, and the LFW dataset. Each image is converted into a protected template and mapped to a facial embedding using the student model. For each embedding, five images are generated using Arc2Face to account for stochasticity. We report both per-image success rate and Success@5. Face generation is highly efficient due to the small backbone of SD1.5: generating a batch of five images for a single embedding takes approximately three seconds on an NVIDIA A6000 GPU, corresponding to a throughput of roughly 1,200 identities per hour and an estimated cost of \$0.0005 per identity generation. Visual examples are shown in Figure~\ref{fig:vis}.



Following the evaluation protocol of CanFG, we employ a commercial face verification system, Face++, to assess identity consistency between the original dataset image and each generated face. We used Face++, marketed to have ``financial-grade security standards," which is usually a higher standard (e.g., more challenging for us) than many open-source methods. This also avoids using the same model (e.g., ArcFace) or models trained on the same datasets (most open source ones) for both embedding extraction and verification. Face++ outputs a confidence score together with three operating thresholds corresponding to error rate of $1\times10^{-3}$, $1\times10^{-4}$, and $1\times10^{-5}$. For each generated image, we record the strictest threshold at which the identity match is accepted and use this as the evaluation outcome. If no face is detected in the original image, it is excluded from the data, while if no face is detected in the generated image, it counts as a failure. 

The results are summarized in Table~\ref{tab:res}. On all three datasets, the success rate at the first attempt is all higher than 97\%, and the success rate for five attempts ranges from 97.9\% to 100\%. Even at the strictest threshold, the success rate is still above 90\%.  We directly compare our regeneration attack with the original reconstruction attack protocol used by FracFace's authors. In Table~\ref{tab:baseline}, we reported the protection rates claimed in FracFace under its own evaluation and the corresponding rates under our attack on the TPDNE dataset.

The evaluation of attack success in FracFace \cite{dai_fracface_2025} is based on a Protection (\%) metric, defined as the proportion of frequency-domain channels that are filtered or structurally disrupted. This formulation establishes a lower barrier for defensive claims than our identity-centric standard, which requires successful regeneration of images passing commercial-grade verification. By our metric, the protection rate of most recent PPFR methods in U-Net/StyleGAN attack is almost always 100\%. Despite our stricter criterion being unfavorable to reported attack success, we show that high channel protection does not prevent identity leakage: even when FracFace claims high protection under its frequency-domain metric, FaceLinkGen achieves near-total identity recovery. This also shows that channel distubution does not mean identity protection. 

To cross-verify this result, we used another commercial facial comparison API from Amazon through EdenAi on 700 selected images on the LFW dataset. The Amazon API only provides a single pass/fail decision with a confidence score; the results are shown in Table~\ref{tab:amazon}. The values are close to the Face++ results, validating our claims. 

To rule out the dependence on models like Arc2Face or third-party verification services like Face++ or Amazon, we compared the similarity of the extracted embeddings with the original face. As detailed in Section~\ref{sec:sim}, the embedding extracted from a protected template shows higher cosine similarity to its source image than to another image of the same person.





\begin{table*}[h!]
    \centering
    % \caption{Comparison with baseline attack methods. Protection rates reported by FracFace under its reconstruction-based evaluation are contrasted with those obtained under our regeneration attack on the TPDNE dataset. While FracFace defines protection failure based on the proportion of preserved frequency-domain channels, our evaluation adopts a stricter criterion: a regenerated face is considered a successful regeneration (failed protection) if it passes Face++ verification at a FAR of 1$\times10^{-3}$. Despite this more conservative definition of attack success, the observed protection rates are substantially lower.}
    \caption{Comparison with FracFace defensive claims. While FracFace \citep{dai_fracface_2025} measures protection by frequency channel disruption, we evaluate actual identity leakage. High disruption rates fail to prevent extraction, as FaceLinkGen achieves near-total recovery through commercial-grade verification.}
    \label{tab:baseline}
    
    \begin{tabular}{cl>{\centering\arraybackslash}p{3cm}>{\centering\arraybackslash}p{3cm}>{\raggedright\arraybackslash}p{3cm}}
    \toprule
         &   Venue&Protection Tested in FracFace \cite{dai_fracface_2025}& Protection Tested Using Our Method &Protection Tested Using Our Method (5 trials)\\
         \midrule
         PartialFace \cite{mi_privacy-preserving_2023} &   ICCV 2023&0.680& 0.002&0.000\\
         MinusFace \cite{mi_privacy-preserving_2024} &   CVPR 2024&0.850& 0.011&0.004\\
         FracFace \cite{dai_fracface_2025}&   NeurIPS 2025&1.000& 0.015&0.008\\
         \bottomrule 
    \end{tabular}
    \vskip -0.1in
\end{table*}


\begin{table}[]
    \centering
        \caption{Regeneration Success@5 on Face++ and Amazon API and 1-to-1 Linkage Success Rate In Assumption-Constrained Settings. Matching accuracy is reported as the max accuracy obtained by optimal threshold.}
    
    \label{tab:blackbox}
    \begin{tabular}{cccc}
    \toprule
      Method &   Face++& Amazon API &Matching\\
         \midrule 
       FracFace&   0.946&  0.473&0.949\\
       PartialFace&   0.946&  0.447&0.925\\
       MinusFace&   0.963&   0.570&0.962\\
             \bottomrule
             \end{tabular}
             %% \vspace{-0.7cm}
               \vskip -0.1in
\end{table}

\section{What If the Attacker Knows Almost Nothing?} \label{ref:bb}

Our main results demonstrate that identity information can be reliably extracted by the primary adversary in the PPFR threat model, namely the service provider or developer who has full access to the conversion process \cite{erkin_privacy-preserving_2009}. This also follows the threat models in prior PPFR evaluations \cite{mi_privacy-preserving_2024,mi_privacy-preserving_2023,dai_fracface_2025}. We additionally ask how far this vulnerability extends to more constrained adversaries: an external attacker or a knowledge-constrained insider who has neither access to the conversion process nor any knowledge of its internals.

To investigate this, we consider an extreme, minimal-assumption scenario. In this setting, the attacker possesses only 30 paired image–template samples for validation (not for training) and has zero knowledge of the underlying protection mechanisms. This is actually stricter than the ``black‑box" scenario in \citet{mi_privacy-preserving_2023}, which assumes the attacker knows the conversion process but not the channel parameters, and more realistic than \citet{zhang_validating_2024}, which relies on thousands of server queries per identity. In real attacks, these 30 pairs can be obtained by a small number of leaked samples, known identities, or attacker-controlled accounts. It can also be simulated with low-frequency queries to the authentication server.

We observe that despite their claimed algorithmic complexity, the output templates of these systems share a common visual essence: they all preserve high-frequency information while obfuscating low-frequency information. Based on this intuition, the attacker can bypass any system-specific modeling and instead use a generic Gaussian-blur-based high-pass filter as a universal proxy task. We avoided using DCT or DWT to decouple from the methods used in the tested systems. By subtracting a slightly blurred version from the original image and applying simple data augmentations (e.g., varying kernels and strengths), the attacker trains a student model to align this simple high-pass domain with the identity embedding space. The attacker does not need to know any details about the system; instead, the high-pass characteristic is easily observable visually.

The training process is the same as our main text, with simply the known PPFR conversion process replaced with a high-pass filter. During inference, the templates are directly fed into the student model except for MinusFace, for which it is passed through a Gaussian-blur-based high-pass filter to remove low-frequency noise.

We trained only one model to attack all three methods. As shown in Table~\ref{tab:blackbox}, identity leakage remains strikingly persistent. For all three systems, we achieved over 92\% 1-to-1 matching success rate and over 94\% re-generation success@5 on Face++ and about 44-57\% on the Amazon API. The Amazon API is likely more strict or sensitive to AI-generated images. Notably, the 1-to-1 linkage success rate on LFW remains around 92-96\% (Table~\ref{tab:blackbox}), close to our main experiments. 


These results indicate that identity-consistent regeneration and reliable linkage remain feasible even under extremely constrained attacker assumptions for frequency-domain PPFR methods. The three evaluated systems, despite their claimed algorithmic complexity, share a common vulnerability: their output representations exhibit strong coupling with simple high-pass filtering operations. We note that this particular proxy-based attack is specific to the frequency-domain family, as the high-pass filter assumption does not apply to non-frequency methods such as CanFG. However, CanFG remains vulnerable under our main oracle-access threat model (Section~5), demonstrating that the vulnerability is not an artifact of frequency-domain design alone.

These results show that, for the frequency-domain PPFR methods evaluated, identity extraction remains feasible even under near-zero-knowledge attacker assumptions. The three systems share a common characteristic: their output representations retain strong coupling with simple high-pass filtering operations, which is what makes this proxy-based attack possible. We note that this particular attack is specific to the frequency-domain family and does not apply to non-frequency methods such as CanFG. However, CanFG remains vulnerable under the oracle-access threat model (Section~5), indicating that identity leakage is not unique to frequency-domain design. Across both settings, the results suggest that as long as templates preserve recognition utility, they tend to retain recoverable identity information, and current frequency-domain defenses do not adequately address this.







    































%
%
\begin{figure}
    \centering
    \includegraphics[width=0.7\linewidth]{imgs/dist.png}
    \caption{Similarity Distribution}
    \label{fig:sim}
\end{figure}

\section{Similarity Distribution}
\label{sec:sim} 
To directly quantify identity leakage from protected templates, independent of the downstream face generation process (Arc2Face) or the specific behavior of commercial verification APIs (Face++, Amazon), we analyze the cosine similarity in the standard ArcFace embedding space: the similarity between two normal images of the same person, and the similarity between one image and its protected template. The histogram is shown in Figure~\ref{fig:sim}. Due to the dataset noise, some real photos and their similarity is near 0 (See \cite{wang_devil_2018}); this also appears in the ArcFace paper \cite{deng_arcface_2022}, but we focus on the main cluster here. We tested this on our testing set. We observed that the similarity between the original image and its template is higher than the similarity between another image of the same person in all three methods. This means that the template is a better identity descriptor for this specific image than another image of the same person. Note that this does not imply that protected templates are universally closer to the underlying identity than real images. Instead, the template remains most similar to its corresponding source image. This indicates that the template is an image-conditioned projection that preserves identity while retaining instance-specific bias, rather than a global identity prototype.



\section{Soft Identity Leakage: Beyond Unique Identifiers}
Beyond hard identity recognition, the exposure of soft biometric attributes presents a significant privacy risk. Characteristics such as skin color, age, and gender are sensitive personal data that enable unauthorized profiling and algorithmic discrimination. Privacy frameworks like the Canadian Privacy Act \citep{branch_consolidated_2025} explicitly protect race and age. A robust privacy-preserving system must therefore prevent the recovery of these attributes from its templates.

In the public rebuttal of FracFace \cite{dai_fracface_2025} on OpenReview \cite{noauthor_wayback_nodate}, the authors claimed successful obfuscation of age, gender, and ethnicity. They cited a human perception study where participants reported less than 13\% usage of these biometrics for identity inference, with over 76\% of participants relying on guesswork. However, our empirical results in Figure~\ref{fig:vis} demonstrate that these soft biometrics remain visible in the re-generated faces. This matches with previous research, which indicates that ArcFace embeddings retain such information \cite{melzi_multi-ive_2023, osorio-roig_reversing_2023}. Because our extracted embeddings closely resemble the original facial embeddings, we hypothesize that a model can learn a direct mapping from the embedding to these attributes without facial reconstruction.

To test this, we trained MLP models on the FairFace dataset \cite{karkkainen_fairface_2019} to predict age, gender, and race (7 classes) across 500 test images. Table~\ref{tab:soft} shows that gender is identified with at least 82\% accuracy, and the Age MAE ranges from 6.1 to 7.5 years. The race accuracy is lower, ranging from 0.50 to 0.60, but considering 7 classes (random baseline: 14.3\%), this represents substantial leakage. We note that the question is not whether ArcFace embeddings contain demographic information in general (they do), but whether the PPFR template retains sufficient information for it to be extracted. If the transformation had genuinely removed soft biometric information, no downstream model could recover it from the protected template, regardless of the model's own biases. The fact that recovery is possible at above-chance accuracy indicates retention in the template, whether explicitly or through correlation with preserved features such as facial geometry. Given that FairFace labels use 10-year intervals, an MAE below 10.0 suggests leakage that matches the inherent precision of the reference model. We also compared this with the ArcFace embeddings computed on the original images as upper bound. As expected, those embeddings perform better than the ones derived from the protected templates, but the performance gap is small enough to indicate that the templates still leak meaningful soft biometric information.

\begin{table}[h]
\centering

\caption{Soft Biometrics Leakage, we included original image (ArcFace) embedding as a baseline}
\label{tab:soft}
\begin{tabular}{lccc}
\toprule
Method & Race Acc $\uparrow$ & Gender Acc $\uparrow$ & Age MAE $\downarrow$ \\
\midrule
FracFace & 0.50& 0.82& 7.5\\
MinusFace & 0.56& 0.86& 6.4\\
PartialFace & 0.60& 0.88& 6.1\\
\midrule
ArcFace & 0.70& 0.94&  4.9\\
\bottomrule
\end{tabular}
\vskip -0.1in
\end{table}

Some methods, such as CanFG \cite{wang_make_2024} and FaceAnonyMixer \cite{alam_faceanonymixer_2025}, intentionally preserve soft biometrics for auxiliary tasks. We contend that this design is problematic. PPFR systems should rely only on identity-discriminative features. Retaining soft biometrics offers attackers more data for reconstruction and profiling without improving recognition performance. These attributes are central to privacy frameworks and require stricter protection \cite{osorio-roig_attack_2022}.

\section{Transfer to Attacking De-Identification Systems}
Our primary focus is on PPFR systems, which retain machine-readable identity information while rendering the original face unreconstructible from the template by removing some visual information. However, a complementary line of work pursues the opposite objective: preserving visual recognizability to humans while disrupting machine-based face recognition \cite{salar_enhancing_2025, yang_towards_2021, shamshad_clip2protect_2023}. These systems apply adversarial noise or synthetic makeup so that the protected face remains human-recognizable but causes FR systems to fail.

This design introduces a vulnerability symmetric to the one we identify in PPFR. PPFR systems leak identity by preserving machine-readable mutual information; these perception-consistent de-identification systems leak identity by preserving human-perceivable information. In both cases, sufficient identity signal is retained for extraction. Notably, the same learning-based pipeline suffices to recover this information in both settings. As long as this mutual information can be aligned to a publicly accessible embedding space, a student model can recover it. 

We selected TIP-IM \cite{yang_towards_2021} as a representative dodging method. We note that many methods \cite{salar_enhancing_2025, hu_protecting_2022} conflate dodging with impersonation, which are not equivalent: an FR system can simultaneously associate a face with both the original identity and the impersonation target. In privacy-purposed de-identification scenarios, users care about dogging andimpersonation, which is closer to attacking an FR system. We evaluate on a small subset of CASIA-WebFace (2,082 images from 408 identities). 

We report $\text{sim}(f(x), f(x_p))$, where $x$ is the original image and $x_p$ is the protected image, reflecting how well the de-identification method deceives the original FR model, and $\text{sim}(f(x), d(x_p))$, measuring whether our distillation pipeline can recover the identity information that the protection attempts to suppress. The results in Table~\ref{tab:deid} show that although TIP-IM provides perfect protection against ArcFace, this is almost completely undermined by our distillation pipeline. This confirms that adversarial perturbation suppresses identity in one model's embedding space while leaving it recoverable by a student aligned to a different space, suggesting a structural limitation of adversarial de-identification as a privacy mechanism. From this perspective, our pipeline in the PPFR domain can be seamlessly viewed as adversarial training \cite{goodfellow_explaining_2015} when targeting perception-consistent de-identification systems. This aligns with well-known security principles: adversarial protection methods will almost always fail once publicly disclosed, and new models can be trained on it \cite{athalye_obfuscated_2018}. It is just a turn-based arm race \cite{biggio_wild_2017}.




\begin{table}[t]
\caption{Identity leakage evaluation on TIP-IM de-identification. We report cosine similarity and linkage accuracy using the standard FR model $f$ and our student model $d$ on protected images $x_p$.}
\label{tab:deid}
\centering
\begin{tabular}{lcc}
\toprule
& $f(x_p)$ & $d(x_p)$ \\
\midrule
Cosine Similarity to $f(x)$ & $-0.081$& $0.648$\\
Linkage Accuracy (same image) & $0.000$& $0.997$ \\
Linkage Accuracy (cross image) & $0.000$& $0.939$\\
\bottomrule
\end{tabular}
\end{table}


\section{Future Directions}
\label{sec:fur}
We suggest several potential pathways for future PPFR designs and evaluations, primarily focusing on stronger defensive mechanisms and broader vulnerability assessments.

\textbf{Cryptographic and Key-Based Hardening.} 
One rigorous approach is to incorporate secret keys into the conversion process, similar to \citet{yuan_generating_2022}. This serves as a multi-factor authentication system (requiring both biometrics and a key), preventing attackers (including our method) from converting a face without the secret. Alternatively, systems may revert to formal cryptographic methods like \cite{ao_cryptoface_2025}. While traditionally viewed as computationally expensive, modern resources make this trade-off acceptable; for instance, \citet{jindal_secure_2020} reports only 2.83ms processing time per face pair. Importantly, these computational costs can act as an effective client-side constraint against brute-force attacks (conceptually similar to slow hashing), enhancing privacy while remaining imperceptible to regular users. Again, we want to emphasize that de-ID or reversible face encryption methods \textit{cannot} be used for PPFR tasks, as they either make the face completely useless for recognition or the reversibility compromises the privacy-preserving nature.


\textbf{De-identification: Fooling Human Perception Instead of Machines.} Current de-identification methods aim to fool FR systems while keeping the identity visually intact to human observers. However, this is almost guaranteed to be attacked due to the information still being intact. However, research suggests that human and machine perception of facial similarity diverges \cite{sadovnik_finding_2018}. We propose an alternative direction: rather than hiding identity from machines while preserving it for humans, one could actively remove identity information from the image while optimizing a human perceptual loss to make the result appear identity-consistent to human observers. This inverts the typical de-identification objective and may offer stronger privacy guaranties if real identity information can be completely removed. One possible approach is to reconstruct the face from identity-agnostic geometric cues such as depth maps or Canny edge maps, which preserve appearance structure without encoding biometric identity. This mapping could also be N-to-1 such that a human can still \textit{recognize} the face, but neither a human nor a machine can re-identify the face due to lost information.






\section{Conclusion}
This paper shows that pixel-level reconstruction metrics do not capture identity-level leakage in frequency-domain PPFR systems. Identity information can be extracted from protected templates via a straightforward distillation pipeline and used for linkage and face regeneration without recovering original pixels. FaceLinkGen achieves over 98.5\% matching accuracy and above 96\% regeneration success on three frequency-domain PPFR systems, suggesting that these methods do not provide the level of identity privacy their pixel-level evaluation results imply.

We hope these findings encourage the adoption of identity-centric evaluation alongside existing metrics in PPFR research, and motivate exploration of stronger protection mechanisms, such as cryptographic approaches, that offer more formal privacy guarantees.
\newpage




 





